\section{Laboratory instrumentation}
\subsection{Oscilloscope}
\subsubsection*{Introduction}
An oscilloscope is an electronic instrument that displays electrical signals graphically as a voltage versus time plot. Some common
use cases of oscilloscopes include:
\begin{itemize}
	\item \textbf{wave shape}: observing the shape of electrical signals
	\item \textbf{amplitude}: measuring the voltage levels of signals
	\item \textbf{frequency}: measuring the frequency of periodic signals
	\item \textbf{time between events}: measuring the time interval between two events in a signal
	\item \textbf{AC/DC components}: analyzing the DC and AC components of signals
	\item \textbf{noise}: detecting and analyzing noise in electrical signals
\end{itemize}

\subsubsection*{Oscilloscope architecture}
\begin{itemize}
	\item \textbf{input coupling}: switch that add a capacitor in series to the input signal to block the DC component and only
	let the AC component pass through, it is often used to isolate the AC signal from the DC offset and to analyze the AC
	characteristics of the signal
	\item \textbf{amplitude control}: amplifier/attenuator that adjusts the amplitude of the input signal to provide a
	suitable voltage range for the analog-to-digital converter (ADC)
	\item \textbf{channel isolation}: optical (with photodiodes) or galvanic (with transformers) isolation to isolate the input
	signal from the internal circuitry of the oscilloscope and to prevent interference and damage to the device
	\item \textbf{Analog to Digital Converter (ADC)}: samples the analog input signal at regular intervals and converts it into
	digital values for processing and display in real time
	\item \textbf{memory}: stores the digital values of the signal for processing and display
	\item \textbf{processor}: manage the sample storage in memory, processes the digital values of the signal to perform various
	functions such as triggering, measurement, graphics processing, \dots and controls the user interface
	\item \textbf{triggering}: digital elaboration of the input signal to detect specific events (rising edge, falling edge,
	edge delay, pulse width, n-cycles, \dots) and align the signal on the screen based on the detected events
\end{itemize}

\subsubsection*{Oscilloscope performances}
\begin{itemize}
	\item \textbf{bandwidth}: the frequency range of the oscilloscope, measured in MHz; the upper limit is the frequency at which
	a sinusoidal signal is attenuated by 3 dB (about 70.7\% of the original amplitude); it is recommended to use an oscilloscope
	with a bandwidth at least 5 times higher than the frequency of the signal to be measured to ensure accurate
	reproduction of the signal 
	\item \textbf{sampling rate}: the number of samples taken per second to reconstruct the input signal; as for bandwidth, it is
	recommended to have a sampling rate at least 5 times higher than the frequency of the signal to be measured
	\item \textbf{record length}: the number of samples that can be stored in memory for analysis and display; it determines the
	maximum time interval that can be observed on the oscilloscope screen at maximum sampling rate, obtained by the formula: 
	\(\textit{acquired time} = \textit{record length} \; / \; \textit{sampling rate}\); if the time window is bigger than the
	acquired time, the oscilloscope will automatically reduce the sampling rate to fit the signal in the time window, which can
	result in a loss of detail and accuracy in the signal representation
\end{itemize}
NOTE: the oscilloscope acquisition is not continuous, it has a dead time between acquisitions during which the oscilloscope
is busy to flush the memory and prepare for the next acquisition.

\subsubsection*{Oscilloscope control panel}
\begin{itemize}
	\item \textbf{channel inputs}: BNC connectors (coaxial cables) for connecting probes to the oscilloscope
	\item \textbf{volts/div}: adjust the vertical scale of the signal on the screen, it determines how many volts are represented
	by each vertical grid division
	\item \textbf{vertical position}: adjust the vertical position of the signal on the screen, it allows to move the signal up
	or down to align it with the grid or to compare it with other signals
	\item \textbf{time/div}: adjust the horizontal scale of the signal on the screen, it determines how much time is represented
	by each horizontal grid division
	\item \textbf{horizontal position}: adjust the horizontal position of the signal on the screen, it allows to move the signal
	left or right to align it with the grid or to compare it with other signals
	\item \textbf{trigger level}: adjust the voltage level at which the oscilloscope triggers and starts the acquisition of the
	signal
	\item other controls for trigger options, coupling, autoscale, cursors, math functions, \dots
\end{itemize}

\subsubsection*{Oscilloscope display}
The oscilloscope display is divided into a grid of horizontal and vertical divisions, which are used to measure the amplitude
and time of the signal. The vertical divisions represent voltage levels, while the horizontal divisions represent time intervals.
There are also indicators for:
\begin{itemize}
	\item \textbf{waveform markers}: horizontal line that indicates the 0 V level on the screen for every channel
	\item \textbf{waveform position}: vertical line that indicates the zero time point on the screen (where the trigger event occurs)
	\item \textbf{cursors}: vertical and horizontal lines that can be moved by the user to measure voltage levels and time intervals
	\item \textbf{div settings}: indicators for the current volts/div and time/div settings for each channel
	\item other indicators for general settings, measurements, information, \dots
\end{itemize}

\subsubsection*{Triggering}
Triggering is a feature of the oscilloscope that allows to horizontally align the signal on the screen by positioning the zero time
point at a specific event in the signal. The trigger can be set to detect different types of trigger events, such as:
\begin{itemize}
	\item \textbf{rising edge}: when the voltage level becomes greater than a specified threshold
	\item \textbf{falling edge}: when the voltage level becomes smaller than a specified threshold
	\item \textbf{edge delay}: after a specified time delay from the rising or falling edge
	\item \textbf{pulse width}: when the signal pulse width is greater or smaller than a specified value
	\item \textbf{n-cycles}: after a specified number of cycles of the signal
	\item \textbf{digital patterns}: recognize typical patterns of  digital protocols (I2C, SPI, CAN, USB, \dots)
\end{itemize}
NOTE: it is not possible to trigger two signals at the same time, if two signals have different frequencies, only one of them will
be triggered and correctly displayed on the screen.

\subsubsection*{Cursors}
Cursors are vertical and horizontal lines that can be moved by the user to measure voltage levels and time intervals on the
oscilloscope display. They can be used to measure amplitude, frequency, period, rise time, fall time, duty cycle, \dots of
the signal. It is a more accurate measurement method compared to counting the number of divisions on the grid.

\subsubsection*{Acquisition mode}
There are different acquisition modes that determine how the oscilloscope samples and processes the input signal. Particularly
important is the \textbf{average mode}, which takes multiple acquisitions of the signal and displays an average of them to filter
high frequency noise and improve the signal quality. Note that the average mode slow down the acquisition process.

\subsubsection*{AC/DC coupling}
The AC/DC coupling switch allows to add a capacitor in series to the input signal to filter out the DC component. This
is useful to analyze the AC characteristics of the signal without being affected by the DC offset. It is important to
note that when the signal frequency is too low, the capacitor starts to filter out low frequencies (10-20 Hz) that are
part of the AC signal, causing voltage ``droop'' or ``sag'' in the waveform. In general it is recommended to use AC coupling
for signals with frequencies over 50-200 Hz.

\subsubsection*{Input impedance and amplification/attenuation stage}
In normal working conditions, the input impedance of the oscilloscope is set to HiZ (high impedance) with a value of
\(1 \; M\Omega\). Some oscilloscopes can also have a low impedance setting (LoZ) with a value of \(50 \; \Omega\) that
bypasses the conditioning stage and directly connects the input signal to the ADC. This setting is useful for measuring
high frequency signals but it can cause damage to the oscilloscope if the input voltage is bigger than 5 V.

\subsection{Function generator}
\subsubsection*{Introduction}
A function generator is an electronic instrument that generates electrical waveforms of adjustable shapes, frequencies and
amplitudes. It acts as an AC voltage source in series with a DC offset voltage source and an internal resistance.

\subsubsection*{Output impedance}
The output impedance of a function generator is \(50 \; \Omega\) and the instrument expects to be connected to a load with
the same impedance. For this reason the real output voltage is doubled to ensure the desired voltage across the load. When
the function generator is connected to a load with a different impedance, the load impedance must be set accordingly. In
particular, when it is connected to a HiZ load, the instrument must be set to HiZ mode.

\subsection{Multimeter}
The multimeter is an electronic instrument that can measure voltage, current and resistance. It is used to measure a wide
voltage/current range with high accuracy. 

\subsection{Power supply}
The power supply is an electronic instrument that provides a stable and adjustable DC voltage to power electronic circuits.
Sometimes it is useful to link the GND terminal to the COM terminal (both of the power supply) to create a common reference
point for the circuit.
